% Part: computability
% Chapter: machines-computations
% Section: configuration

\documentclass[../../../include/open-logic-section]{subfiles}

\begin{document}

\olfileid[pt]{cmp}{tur}{con}
\olsection{Configurações e Computações}

\begin{explain}
Lembre-se de acompanhar as configurações da máquina dos pares
anteriormente. O mecanismo imaginário formado pela fita, pela cabeça de
leitura-escrita e pelo programa da máquina de Turing é, na verdade, apenas
uma maneira intuitiva de visualizar o que é uma computação de máquina de
Turing. Formalmente, podemos definir a computação de uma máquina de Turing
sobre uma dada entrada como uma sequência de
\emph{configurações}---e uma configuração, por sua vez, é uma sequência de
símbolos (correspondendo ao conteúdo da fita em um dado ponto da
computação), um número que indica a posição da cabeça de leitura-escrita
e um estado. Com isso, podemos definir o que a máquina de
Turing~$M$ computa sobre uma dada entrada.
\end{explain}

\begin{defn}[Configuração]
Uma \emph{configuração} da máquina de Turing $M = \tuple{Q, \Sigma, q_0,
\delta}$ é uma tripla $\tuple{C, m, q}$ onde
\begin{enumerate}
\item $C \in \Sigma^*$ é uma sequência finita de símbolos de $\Sigma$,
\item $m \in \Nat$ é um número $< \len{C}$, e
\item $q \in Q$
\end{enumerate}
Intuitivamente, a sequência~$C$ é o conteúdo da fita (símbolos de
todos os quadrados desde o quadrado mais à esquerda até o último não
vazio ou previamente visitado), $m$~é o número do quadrado que a
cabeça de leitura-escrita está examinando (começando com $0$ como o número do
quadrado mais à esquerda), e $q$ é o estado atual da máquina.
\end{defn}

\begin{explain}
A possível entrada para uma máquina de Turing é uma sequência de símbolos,
geralmente uma sequência que codifica um número de alguma forma. A
configuração inicial da máquina de Turing é aquela configuração em que
iniciamos a máquina de Turing para trabalhar sobre essa entrada: a fita contém o
marcador de fim de fita imediatamente seguido da entrada escrita nos
quadrados à direita, a cabeça de leitura-escrita está examinando o quadrado
mais à esquerda da entrada (ou seja, o quadrado à direita do marcador de
fim esquerdo), e o mecanismo está no estado inicial designado~$q_0$.
\end{explain}

\begin{defn}[Configuração inicial]
A \emph{configuração inicial} de $M$ para a entrada $I \in \Sigma^*$ é
\[
\tuple{\TMendtape \frown I, 1, q_0}.
\]
\end{defn}

\begin{explain}
O símbolo~$\frown$ denota \emph{concatenação}---a cadeia de entrada
começa imediatamente após o marcador de fim esquerdo.
\end{explain}

\begin{defn}
Dizemos que uma configuração $\tuple{C, m, q}$ \emph{produz a
configuração $\tuple{C', m', q'}$ em um passo} (de acordo com~$M$),
sse
\begin{enumerate}
\item o $m$-ésimo símbolo de $C$ é $\sigma$,
\item o conjunto de instruções de $M$ especifica $\delta(q, \sigma) =
  \tuple{q', \sigma', D}$,
\item o $m$-ésimo símbolo de $C'$ é $\sigma'$, e 
\item
\begin{enumerate}
\item $D = L$ e $m' = m - 1$ se $m>0$, caso contrário $m'=0$, ou
\item $D = R$ e $m' = m + 1$, ou
\item $D = N$ e $m' = m$,
\end{enumerate}
\item se $m' = \len{C}$, então $\len{C'} = \len{C} + 1$ e o $m'$-ésimo
  símbolo de $C'$ é~$\TMblank$. Caso contrário $\len{C'}=\len{C}$.
\item para todo $i$ tal que $i < \len{C}$ e $i \neq m$, $C'(i) = C(i)$,
\end{enumerate}
\end{defn}

\begin{defn}\ollabel{defn:run-output}
Uma \emph{execução de $M$ sobre a entrada~$I$} é uma sequência $C_i$ de configurações
de $M$, onde $C_0$ é a configuração inicial de $M$ para a entrada~$I$,
e cada $C_i$ produz $C_{i+1}$ em um passo.

Dizemos que $M$ \emph{para sobre a entrada $I$ após $k$ passos} se $C_k =
\tuple{C, m, q}$, o $m$-ésimo símbolo de~$C$ é~$\sigma$, e $\delta(q,
\sigma)$ é indefinido. Nesse caso, a \emph{saída} de~$M$ para a
entrada~$I$ é~$O$, onde $O$ é uma cadeia de símbolos que não termina
em~$\TMblank$ tal que $C = \TMendtape \concat O \concat \TMblank^j$
para algum~$j \in \Nat$. ($\TMblank^j$ é uma sequência de $j$ $\TMblank$'s.)
\end{defn}

\begin{explain}
De acordo com esta definição, a saída~$O$ de~$M$ sempre 
termina em um símbolo diferente de~$\TMblank$, ou, se no instante~$k$ toda a
fita estiver preenchida com~$\TMblank$ (exceto o~$\TMendtape$ mais à esquerda),
$O$~é a cadeia vazia.
\end{explain}

\end{document}
